\section{Introduction}
\label{sec-intro}
Nowadays, recommender systems play an indispensable role in meeting users’ personalized interests and alleviating the information overload problem.
In large-scale industrial recommendation, the recommender model usually has to deal with multiple scenarios (\eg homepage feeds or banner) and multiple user behaviors (\eg click or like), referred to as multi-scenario learning (MSL)~\cite{msl-benchmark} and multi-task learning (MTL)~\cite{mtl-survey} in recommender system.

In recommendations, multi-scenario learning (MSL)~\cite{PLE,HMoE,STAR,PEPNet,hinet,mmfi,m3oe,ADL} aims to learn data from multiple scenarios within a unified model, and multi-task learning (MTL)~\cite{sharedbottom,MMoE,Cross-Stitch,snr,adatt} aims to optimize the unified model with multiple related objectives.
Technically speaking, the key of MSL and MTL in recommendations lies in the modeling of multi-distribution interrelations, capturing the commonalities as well as discriminating the differences across scenarios and tasks.
Most existing methods usually adopt a \emph{shared-specific} framework, where scenario/task-shared and scenario/task-specific components are combined in the backbone.
For example, one class of methods~\cite{MMoE,HMoE,PLE} is based on the Mixture-of-Experts (MoE) structure, which models different scenarios and tasks by designing both shared experts and scenario/task-specific expert.
Another class of methods~\cite{STAR,PEPNet} is based on the dynamic parameter generation approach, where scenario- or task-specific prior information is used to generate dedicated network parameters (such as MLPs or Gate units), enabling differentiated modeling across scenarios and tasks.

Although existing methods have achieved notable results, there are still several issues that remain to be addressed.
(1) \emph{Limitation in scaling capability}: inspired by the scaling law in large language models (LLMs)~\cite{LLM1,LLM2}, recent industrial recommender systems focus on improving model learning capability by scaling up the model parameters of feature-interaction modules, achieving remarkable results~\cite{RankMixer,KGBRank,OneTrans,wukong,HSTU,MTGR,hyformer}.
Typically, the parameter size of the state-of-the-art models is now hundreds of times larger than that of traditional recommender models, usually exceeding 0.1B parameters~\cite{RankMixer,KGBRank}.
Based on large-parameter recommender models, existing MSL and MTL methods lack effective utilization of complex feature interaction modules.
Because these methods typically only affect certain intermediate modules or shallow output layers, making it difficult to fully leverage the benefits of larger model parameters.
(2) \emph{Difficulty in uniformity modeling}: multi-scenario modeling primarily focuses on differences in input distributions, while multi-task modeling emphasizes differences in label distributions.
Traditional shared-specific structures are often designed differently for these two aspects.
It is challenging to jointly model scenario information and task information in a unified way.

To address the above issues, we further draw inspiration from LLMs to enhance MSL and MTL in recommendations.
LLMs accomplish specific downstream tasks by inputting appropriate prompt information~\cite{gpt4,qwen3}.
As a special type of token, prompt tokens can fully activate the vast latent knowledge of LLMs through the self-attention mechanism.
Similarly, we posit that the scenario and task information can serve as prompts within recommender models like that in LLMs.
By tokenizing scenario and task priors into a unified latent space, we transform them from mere auxiliary inputs into active participants within the core feature-interaction backbone.
Instead of restricting scenario/task information to the "shallows" of the model (e.g., gates or heads), tokenized priors can penetrate deep and multi-layered interaction modules. This allows distribution-specific signals to adaptively modulate the model's massive parameter space, fully harvesting the benefits of the scaling law (\emph{corresponding to the first issue}).
By projecting both input-side scenario contexts and output-side task objectives into a consistent tokenized format, we bridge the gap between MSL and MTL. This enables a holistic and joint modeling framework where scenario-task dependencies are captured through a unified interaction protocol (\emph{corresponding to the second issue}).

To this end, we propose the \textbf{M}ulti-\textbf{D}istribution \textbf{L}earning (\textbf{MDL}) framework, a unified approach for multi-scenario and multi-task learning in large-scale reommender system.
At its core, MDL adopts a "Tokenize-and-Interact" philosophy: it tokenizes feature, scenario, and task information in a unified way, then elaborate effective interactions among these tokens to activate the model capability for MSL and MTL.
Specifically, we first use a \emph{unified information tokenization} module to resolve the representational gap between disparate distributions.
This module applies semantic segmentation and nonlinear transformations to raw input features, projecting them into a discrete set of feature tokens, scenario tokens, and task tokens.
Based on these tokens, we facilitate multi-granular information exchange through three specialized mechanisms.
The feature token self-interaction adopts a self-attention-like mechanism to capture high-order feature dependencies, providing the recommender with fundamental expressive capabilities.
The interaction between feature tokens and scenario/task tokens is achieved through \emph{domain-aware attention}, 
where scenario and task tokens act as "queries" that adaptively activate and aggregate relevant feature information. 
By doing so, the model learns to prioritize different feature sub-spaces based on the specific scenario/task context, effectively handling distribution shifts.
Finally, the interaction between scenario and task tokens is achieved through the \emph{domain-fused module}, which aggregates token information to enable a unified and robust output for multi-distribution prediction.
By stacking multiple layers, MDL enables feature, scenario, and task information to fully interact in a bottom-up and layer-wise manner.
This deep integration allows the model's massive parameter space to be dynamically steered by distribution-specific prompts, resulting in an effective modeling of complex multi-distribution data.

To evaluate the proposed approach \modelname, we conducted extensive experiments on real-world industrial dataset.
The experimental results demonstrate that our approach can achieve much better performance than competitors for multi-scenario and multi-task recommendations in various settings.
Furthermore, we conducted an online A/B test on Douyin Search for one month, achieving a significant improvement of +0.0626\% in LT30 and -0.3267\% in change query rate.
And the proposed approach MDL has been fully deployed on Douyin Search.
